


\begin{definition}[Divergence]
\citep[pp]{marsden1994mathematical} 
Given a vector field \(\bm{X}\) and a volume form \(\mu\), define the divergence \(\operatorname{div}_{\mu}\bm{X}\) 
\[
(\operatorname{div}_{\mu}\bm{X})\mu \triangleq \mathcal{L}_{\bm{X}}\mu
\equiv \dd (\interior{\bm{X}}{\mu})
\]
\end{definition}

where \(\interior{.}{.}\) is the \emph{interior product} . 

Also 
$\langle\mathbf{x} \wedge \mathbf{v}, \mathbf{w}\rangle=\left\langle\mathbf{v}, \iota_{\mathbf{x}^{\flat}} \mathbf{w}\right\rangle$
And
$\iota_v \omega=*\left(v^{\flat} \wedge * \omega\right)$

\begin{proposition}
Define a new volume form \(\mu_1 = \alpha \mu_2\) for some positive scalar field \(\alpha\). 
Then
\[
\alpha \operatorname{div}_1 \bm{X} = \operatorname{div}_2(\alpha \bm{X})
\]
\[
\alpha \operatorname{div}_{\alpha \mu} \bm{X} = \operatorname{div}_{\mu}(\alpha \bm{X})
\]
and
\[
\operatorname{div}_{\mu} \bm{X} = \operatorname{div}_{(\alpha \mu)} \bm{X}+\bm{X}[\log \alpha]
\]
\end{proposition}
\begin{proof}
% Recall $\mathcal{L}_X\left(\alpha \mu_2\right)= X[\alpha] \mu_2+\alpha \mathcal{L}_X \mu_2$
$$
\begin{aligned}
\left(\operatorname{div}_1 \bm{X}\right) \mu_1 & =\dd \left(\interior{\bm{X}}{} \mu_1\right) \\
& =\dd \left(\interior{\bm{X}}{}\left(\alpha \mu_2\right)\right) \\
& =\dd \left(\alpha \interior{\bm{X}}{} \mu_2\right) \\
& =\dd \alpha \wedge \interior{\bm{X}}{} \mu_2 + \alpha \dd\left(\interior{\bm{X}}{} \mu_2\right) \\
& = \bm{X}[\alpha] \mu_2+\alpha\left(\operatorname{div}_2 \bm{X}\right) \mu_2 \\
& =\left(\frac{1}{\alpha}\bm{X}[\alpha]+\operatorname{div}_2 \bm{X}\right) \mu_1
\end{aligned}
$$
Since $\dd \alpha \wedge \mu_2$ is an ( $n+1$ )-form on an $n$-manifold, it is zero, so

$$
0=\interior{\bm{X}}{}\left(\dd \alpha \wedge \mu_2\right)=\left(\interior{\bm{X}}{} \dd \alpha\right) \wedge \mu_2- \dd \alpha \wedge\left(\interior{\bm{X}}{} \mu_2\right)
$$
Note that \(X[\alpha] = \alpha X[\log \alpha]\) is a consequence of the chain rule \(\nabla_X(g(h)) = g' \nabla_X h\):
\[
X[\log \alpha]
:=\dd (\log \alpha)(X)
= \nabla_X(\log \alpha) 
=\frac{1}{\alpha} \nabla_X \alpha
=\frac{X[\alpha]}{\alpha}
\]

\end{proof}
\iffalse
The divergence operator can then be defined relative to measures \(\volG\) and \(\volE\). 
That is, for a vector field $V$, define the divergence with respect to $\mu$ as:
$$
\operatorname{div}_\mu \boldsymbol{V}:=\frac{1}{n_E} \nabla \cdot\left(n_E \boldsymbol{V}\right)
$$

The elliptic operator associated with a measure $\dd \mu = w\, \dd V$ is precisely:

$$
\operatorname{div}_\mu(\nabla f)=\frac{1}{w} \nabla \cdot(w \nabla f) .
$$
\fi 



\begin{theorem}[Divergence theorem]
Let $M$ be a Riemannian manifold and let $X \in$ $\Gamma_{C^1}(T M)$. Then,
\begin{equation}
\int_{\PlaneManifold} \operatorname{div}_\mu X \, \mu = \int_{\PlaneBoundary} \iota_X \, \mu
\label{eq:gauss-div-thm}
\end{equation}
$$
\int_{\PlaneManifold} \operatorname{div}_g X \omega_g=\int_{\PlaneBoundary}\langle X, \nu\rangle \sigma_g
$$

where $\nu$ is the unit vector normal to $\partial M$.
\end{theorem}
